% Part: sets-functions-relations
% Chapter: sets
% Section: subsets

\documentclass[../../../include/open-logic-section]{subfiles}

\begin{document}

\olfileid{sfr}{set}{sub}
\olsection{ఉపసమితులు, ఘాత సమితులు}

\begin{explain}
సమితులను పోల్చవలసిన సందర్భాలు తరచూ వస్తాయి. ఒక సహజమైన పోలిక
ఇలా ఉంటుంది: \emph{ఒక సమితిలో ఉన్న ప్రతి వస్తువూ మరొక సమితిలోనూ
ఉండటం}. ఈ పరిస్థితి ముఖ్యమైనది కనుక దానిని సూచించడానికి కొత్త
సంకేతాన్ని పరిచయం చేస్తాం.
\end{explain}

\begin{defn}[ఉపసమితి]
$A$ అనే సమితిలోని ప్రతి \tetoken{మూలకం}{!!{element}}, $B$లో కూడా \tetoken{ఒక మూలకం}{!!a{element}}
అయితే, $A$ను $B$ యొక్క \emph{ఉపసమితి} అని చెప్పి
$A \subseteq B$ అని రాస్తాం. $A$ అనేది $B$కు ఉపసమితి కాకపోతే
$A \not\subseteq B$ అని రాస్తాం.
$A \subseteq B$ అయి, $A \neq B$ అయితే $A \subsetneq B$ అని రాసి,
$A$ను $B$ యొక్క \emph{క్రమ ఉపసమితి} అంటాం.
\end{defn}

\begin{ex}
ప్రతి సమితీ తనకు తానే ఉపసమితి. అలాగే $\emptyset$ ప్రతి సమితికీ
ఉపసమితి. సరి సహజ సంఖ్యల సమితి సహజ సంఖ్యల సమితికి ఉపసమితి.
అలాగే $\{ a, b \} \subseteq \{ a, b, c \}$.
కానీ $\{ a, b, e \}$ అనేది $\{ a, b, c \}$కు ఉపసమితి కాదు.
\end{ex}

\begin{ex}
$2$ అనే సంఖ్య పూర్ణ సంఖ్యల సమితిలో \tetoken{మూలకం}{!!{element}}; సరి సంఖ్యల
సమితి మాత్రం పూర్ణ సంఖ్యల సమితికి ఉపసమితి. అయితే ఒక సమితి
మరొక సమితిలో \tetoken{ఒక మూలకం}{!!a{element}} కావడమూ, దానికి ఉపసమితి కావడమూ
\emph{రెండూ} జరగవచ్చు. ఉదాహరణకు $\{0\} \in \{0, \{0\}\}$;
అలాగే $\{0\} \subseteq \{0, \{0\}\}$.
\end{ex}

సమితులు ఒకటే కావడానికి సమితుల సమానత్వ సూత్రం ఒక ప్రమాణాన్ని ఇస్తుంది:
$A = B$ కావడానికి అవసరమైన మరియు సరిపడే షరతు ఏమిటంటే,
$A$లోని ప్రతి \tetoken{మూలకం}{!!{element}}, $B$లో \tetoken{ఒక మూలకం}{!!a{element}} కావాలి;
తిరుగుదిశలోనూ ఇదే వర్తించాలి. ``ఉపసమితి'' నిర్వచనంలోని
$A \subseteq B$ అనేది ఈ ప్రమాణంలోని మొదటి భాగమే:
$A$లోని ప్రతి \tetoken{మూలకం}{!!{element}}, $B$లో \tetoken{ఒక మూలకం}{!!a{element}} కావడం.
$A$, $B$లను పరస్పరం మార్చినా నిర్వచనం వర్తిస్తుంది:
$B \subseteq A$ కావడానికి అవసరమైన మరియు సరిపడే షరతు,
$B$లోని ప్రతి \tetoken{మూలకం}{!!{element}}, $A$లో \tetoken{ఒక మూలకం}{!!a{element}} కావడం.
ఇదే సమితుల సమానత్వ సూత్రంలోని ``తిరుగుదిశ'' భాగం.
అందువల్ల సమితులు ఒకదానికొకటి ఉపసమితులైనప్పుడే సమానమవుతాయి.

\begin{prop}
$A = B$ కావడానికి అవసరమైన మరియు సరిపడే షరతులు
$A \subseteq B$, $B \subseteq A$ రెండూ నెరవేరడం.
\end{prop}

ఇక్కడ మరికొన్ని ఉపయోగకరమైన సంకేతాలను పరిచయం చేద్దాం.
$A$ అనేది $B$కు ఉపసమితి ఎప్పుడవుతుందో నిర్వచించేటప్పుడు,
``$A$లోని ప్రతి \tetoken{మూలకం}{!!{element}} \dots'' అన్న వాక్యంలోని
``$\dots$'' స్థానంలో ``$B$లో \tetoken{ఒక మూలకం}{!!a{element}}'' అని ఉంచాం.
ఈ \emph{రూపం} గల వాక్యాలు తరచూ వస్తాయి కనుక వాటికి
ఒక నియత సంకేతాన్ని ప్రవేశపెట్టడం ఉపయోగకరం.

\begin{defn}\ollabel{forallxina}
$(\forall x \in A)\phi$ అనేది $\forall x(x \in A \lif
\phi)$కు సంక్షిప్త రూపం. అదే విధంగా $(\exists x \in A)\phi$
అనేది $\exists x(x \in A \land \phi)$కు సంక్షిప్త రూపం.
\end{defn}

ఈ సంకేతంతో, $A \subseteq B$ కావడానికి అవసరమైన మరియు సరిపడే
షరతు $(\forall x \in A)x \in B$ అని చెప్పవచ్చు.

ఇప్పుడు ఇచ్చిన సమితి యొక్క ఉపసమితులన్నింటితో ఏర్పడే సమితిని పరిశీలిద్దాం.

\begin{defn}[ఘాత సమితి]
$A$ అనే సమితికి గల ఉపసమితులన్నింటితో ఏర్పడే సమితిని
$A$ యొక్క \emph{ఘాత సమితి} అని పిలిచి, $\Pow{A}$ అని రాస్తాం.
  \[
    \Pow{A} = \Setabs{B}{B \subseteq A}
  \]
\end{defn}

\begin{ex}
$\{ a, b, c \}$కు గల ఉపసమితులన్నీ ఏవి? అవి:
$\emptyset$, $\{a \}$, $\{b\}$, $\{c\}$, $\{a, b\}$, $\{a, c\}$,
$\{b, c\}$, $\{a, b, c\}$. వీటన్నింటి సమితి
$\Pow{\{a,b,c\}}$:
\[
\Pow{\{ a, b, c \}} = \{\emptyset, \{a \}, \{b\}, \{c\}, \{a, b\},
\{b, c\}, \{a, c\}, \{a, b, c\}\}
\]
\end{ex}

\begin{prob}
$\{a, b, c, d\}$కు గల ఉపసమితులన్నింటినీ జాబితాగా రాయండి.
\end{prob}

\begin{prob}
$A$లో $n$ \tetoken{మూలకాలు}{!!{element}s} ఉంటే, $\Pow{A}$లో $2^n$ \tetoken{మూలకాలు}{!!{element}s}
ఉంటాయని చూపండి.
\end{prob}

\end{document}
